%==============================================================================
%  CAPÍTULO XIV
%==============================================================================
\chapter{La entrevista inicial y el diagnóstico del cliente}
\label{cap:entrevista}

\fichaCapitulo{La primera hora define el caso.}{Cuestionario de admisión, verificación
documental independiente, calificación del deudor y elección de la vía procedimental.}

%==============================================================================
\section{Por qué la entrevista inicial decide el resultado}
%==============================================================================

Los procedimientos concursales se pierden, casi siempre, antes de presentarse. No se
pierden en la audiencia inicial ni en la junta de acreedores: se pierden en la primera
reunión, cuando el abogado acepta como verdadero el relato del cliente y construye sobre
él una estrategia que el expediente terminará desmintiendo.

Esta afirmación no es retórica. El deudor en crisis patrimonial llega a la entrevista con
tres distorsiones sistemáticas, y las tres son de buena fe:

\begin{enumerate}
  \item \textbf{Subestima el pasivo.} No miente: ha desplazado de su atención las
  obligaciones que dejó de servir hace tiempo. Las deudas repactadas, las que fueron
  cedidas a empresas de cobranza y las que quedaron con codeudor tienden a desaparecer de
  su relato.

  \item \textbf{Sobreestima el activo.} Valora sus bienes al precio que pagó por ellos o
  al que cree que valen, no al de realización forzada. La maquinaria que compró en 3.000
  UF no se remata en 3.000 UF.

  \item \textbf{Omite los actos del bienio anterior.} No porque quiera ocultarlos, sino
  porque no percibe que sean jurídicamente relevantes. La transferencia del auto a nombre
  del cónyuge, el pago anticipado al proveedor amigo o la hipoteca constituida el año
  pasado sobre una deuda vieja son, para él, hechos domésticos.
\end{enumerate}

Cada una de estas tres distorsiones tiene una consecuencia procesal grave y específica.
La primera destruye la propuesta de acuerdo cuando emergen acreedores no considerados en
el mapa de mayorías. La segunda genera expectativas que ningún procedimiento puede
satisfacer, y con ellas la insatisfacción del cliente y el riesgo de reclamo
profesional. La tercera es, con diferencia, la más peligrosa: expone al deudor a acciones
revocatorias (capítulo \ref{cap:revocatorias}) y, desde la \Lreforma{}, al incidente de
mala fe del \art{169 A} (capítulo \ref{cap:mala-fe}), cuya consecuencia es la denegación
del \discharge{} — es decir, la pérdida completa del beneficio que motivó el encargo.

\begin{foco}[La regla de oro de la entrevista concursal]
Nada de lo que el cliente afirme sobre su patrimonio ingresa al expediente sin respaldo
documental obtenido por el abogado. No por desconfianza, sino porque la declaración
jurada la firma el deudor y la responsabilidad de su exactitud lo alcanza a él, mientras
que la responsabilidad de haberla construido sin verificación es del profesional.
\end{foco}

%==============================================================================
\section{Estructura de la primera reunión}
%==============================================================================

La entrevista inicial no debe intentar resolver el caso. Debe conseguir tres cosas, en
este orden, y la inversión del orden es una fuente frecuente de error.

\subsection{Primer objetivo: calificar al sujeto}

Antes de discutir estrategia alguna es necesario determinar si se está frente a una
\empresadeudora{} o a una \personadeudora{}, y en el primer caso, ante qué categoría de
tamaño. De esa calificación depende íntegramente el catálogo de procedimientos
disponibles (capítulo \ref{cap:estatuto-deudores}). Discutir una reorganización con quien
sólo puede acceder a la renegociación administrativa es perder la reunión.

\subsection{Segundo objetivo: calificar la crisis}

Determinado el sujeto, corresponde establecer si la crisis es de liquidez o estructural
(capítulo \ref{cap:evolucion}, sección 1.3). La distinción no es académica: define si
existe algo que reorganizar o si el encargo consiste en administrar ordenadamente una
liquidación.

\subsection{Tercer objetivo: detectar contingencias}

Sólo entonces corresponde indagar los hechos que pueden frustrar el procedimiento:
actos revocables, activos omitidos, deudas de tratamiento especial y ejecuciones en
curso. Esta indagación se hace al final, cuando el cliente ya ha entrado en confianza y
comprende que la información le sirve a él.

\begin{practica}[Sobre el orden de las preguntas]
Preguntar por los actos del bienio anterior al comienzo de la reunión produce respuestas
defensivas. Preguntarlo después de haber explicado qué es una acción revocatoria y qué
consecuencia tiene ocultar un acto produce respuestas útiles. El cliente colabora cuando
entiende que la pregunta lo protege.
\end{practica}

%==============================================================================
\bloqueEntrevista
%==============================================================================

Lo que sigue es un cuestionario de admisión completo. No todas las preguntas aplican a
todos los casos, pero todas deben formularse: la que parece irrelevante es
frecuentemente la que revela el problema.

\subsection{Bloque A — Calificación subjetiva del deudor}

\begin{cuestionario}
  \pregunta{¿Ha tenido usted o su empresa inicio de actividades en primera categoría en
  los últimos veinticuatro meses?}
  {Determina si es \empresadeudora{} o \personadeudora{}, y con ello todo el catálogo
  procedimental. No se responde con la memoria del cliente: se responde con la carpeta
  tributaria.}

  \pregunta{¿Emite usted boletas de honorarios?}
  {Tras la \Lreforma{} los contribuyentes de segunda categoría acceden a la renegociación
  administrativa. Antes quedaban en tierra de nadie.}

  \pregunta{¿Cuál fue el monto de ventas del último ejercicio comercial?}
  {Define el umbral de la \Lpyme{} y, por tanto, el acceso a los procedimientos
  simplificados, que son sustancialmente más baratos y rápidos.}

  \pregunta{¿De dónde proviene efectivamente su sustento: remuneración dependiente,
  honorarios, pensión, o utilidades de la empresa?}
  {Un inicio de actividades sin movimiento real no necesariamente priva de la calidad de
  \personadeudora{} (capítulo \ref{cap:estatuto-deudores}). Esta pregunta prepara esa
  alegación.}

  \pregunta{¿Existen otras sociedades en las que usted participe o administre?}
  {Detecta grupos empresariales, acreedores relacionados y eventuales confusiones
  patrimoniales que pueden aflorar durante el procedimiento.}
\end{cuestionario}

\subsection{Bloque B — Composición y magnitud del pasivo}

\begin{cuestionario}
  \pregunta{Enumere todas sus deudas, con acreedor, monto aproximado y antigüedad del
  incumplimiento.}
  {Es la línea base que se contrastará contra el informe consolidado. La diferencia entre
  lo declarado y lo documentado es, en sí misma, el dato más útil de la entrevista.}

  \pregunta{¿Alguna de sus deudas fue repactada, renegociada o cedida a una empresa de
  cobranza?}
  {Son las que el cliente sistemáticamente omite. Cambian de acreedor y de RUT, y
  reaparecen en el procedimiento como verificaciones inesperadas.}

  \pregunta{¿Hay deudas en que un tercero figure como codeudor solidario, aval o fiador?}
  {El concurso del deudor principal no libera al codeudor. Si el codeudor es un familiar,
  esta pregunta define la viabilidad emocional del encargo, no sólo la jurídica.}

  \pregunta{¿Tiene deudas previsionales o con la Tesorería General de la República?}
  {Gozan de preferencia y su tratamiento difiere del pasivo común. Las cotizaciones
  impagas, además, tienen consecuencias laborales autónomas.}

  \pregunta{¿Tiene crédito universitario, CAE o Fondo Solidario?}
  {Su suerte frente al \discharge{} es objeto de jurisprudencia contradictoria
  (capítulo \ref{cap:discharge-cae}). El cliente debe saber, desde la primera reunión,
  que probablemente sobrevivirá al procedimiento.}

  \pregunta{¿Tiene deudas por pensión de alimentos?}
  {Requieren advertencia expresa sobre su tratamiento diferenciado y sobre los registros
  especiales de deudores.}
\end{cuestionario}

\begin{alerta}[Las dos preguntas que nunca deben omitirse]
El CAE y las pensiones de alimentos. Son las dos deudas que el cliente da por
extinguibles y que con mayor probabilidad sobrevivirán al procedimiento. Descubrirlo
después de la resolución de término convierte un encargo exitoso en un reclamo.
\end{alerta}

\subsection{Bloque C — Activo y garantías}

\begin{cuestionario}
  \pregunta{Enumere sus bienes: inmuebles, vehículos, maquinaria, existencias, cuentas
  bancarias, instrumentos de ahorro.}
  {La omisión de cualquiera de ellos en el inventario configura la causal del
  \art{169 A} \Nro 1.}

  \pregunta{¿Tiene ahorro previsional voluntario, cuenta 2, depósitos a plazo o fondos
  mutuos?}
  {Es la omisión más frecuente y la más citada por la jurisprudencia como indicio de mala
  fe (capítulo \ref{cap:liquidacion-simplificada}). El cliente no los considera «bienes».}

  \pregunta{¿Existen hipotecas, prendas o prohibiciones sobre alguno de estos bienes?}
  {Los acreedores con garantía real tienen una posición distinta en el concurso y su
  comportamiento condiciona la viabilidad de cualquier reorganización.}

  \pregunta{Si hay bienes gravados: ¿son esenciales para la continuidad del giro?}
  {Determina si una reorganización es siquiera concebible. Una empresa cuyo activo
  productivo está íntegramente gravado rara vez puede reorganizarse.}

  \pregunta{¿Cuál es su remuneración líquida mensual y qué descuentos voluntarios tiene?}
  {En liquidación de persona natural define el tramo inembargable y el excedente
  realizable (capítulo \ref{cap:liquidacion-simplificada}).}

  \pregunta{¿Hay bienes a nombre de terceros que en los hechos sean suyos, o bienes suyos
  que use un tercero?}
  {Pregunta incómoda e indispensable. Anticipa tanto acciones revocatorias como
  tercerías.}
\end{cuestionario}

\subsection{Bloque D — Actos del período sospechoso}

\begin{practica}[Cómo introducir este bloque]
Conviene precederlo de una explicación breve: «la ley permite dejar sin efecto ciertos
actos celebrados antes del procedimiento, y ocultarlos puede costarle el perdón de las
deudas. Necesito conocerlos para protegerlo, no para juzgarlo.»
\end{practica}

\begin{cuestionario}
  \pregunta{En los últimos dos años, ¿vendió, donó o transfirió algún bien?}
  {Período sospechoso de la acción revocatoria objetiva para actos gratuitos, y de la
  subjetiva en general (capítulo \ref{cap:revocatorias}).}

  \pregunta{¿Alguna de esas operaciones fue con familiares, socios o empresas
  relacionadas?}
  {Extiende el período sospechoso y agrava la presunción de perjuicio a la masa.}

  \pregunta{En el último año, ¿pagó alguna deuda antes de su vencimiento, o pagó a un
  acreedor de forma distinta a la pactada?}
  {Pagos anticipados y daciones en pago inusuales son supuestos típicos de revocación
  objetiva.}

  \pregunta{¿Constituyó hipoteca o prenda sobre una deuda que ya existía sin garantía?}
  {Supuesto clásico de revocación objetiva: mejora la posición de un acreedor en
  perjuicio de los demás.}

  \pregunta{¿Ha destruido, extraviado o dejado de llevar documentación contable?}
  {Causal autónoma del \art{169 A}. La respuesta afirmativa exige reconstruir la
  documentación antes de presentar.}
\end{cuestionario}

\subsection{Bloque E — Situación procesal y urgencias}

\begin{cuestionario}
  \pregunta{¿Ha sido notificado de alguna demanda ejecutiva? ¿Cuándo?}
  {En renegociación administrativa incide directamente en la admisibilidad (capítulo
  \ref{cap:renegociacion}). En general, define la urgencia real del encargo.}

  \pregunta{¿Hay embargos trabados, remates fijados o retenciones sobre su
  remuneración?}
  {Determina el calendario. Un remate con fecha convierte un encargo ordinario en uno
  urgente.}

  \pregunta{¿Algún acreedor le ha anunciado una solicitud de liquidación forzosa?}
  {Cambia la estrategia por completo: puede convenir anticiparse con una presentación
  voluntaria.}

  \pregunta{¿Tiene trabajadores con remuneraciones o cotizaciones impagas?}
  {Créditos con preferencia de primera clase y efectos laborales autónomos (capítulo
  \ref{cap:laboral-penal}).}
\end{cuestionario}

\subsection{Bloque F — Objetivos y expectativas del cliente}

\begin{cuestionario}
  \pregunta{¿Qué espera obtener de este procedimiento?}
  {La respuesta revela si la expectativa es realizable. «Que me dejen tranquilo» y
  «que no pierda la casa» exigen respuestas técnicas muy distintas.}

  \pregunta{¿Hay algún bien que necesite conservar por sobre los demás?}
  {Permite anticipar el conflicto entre la expectativa del cliente y el efecto del
  \desasimiento{}.}

  \pregunta{¿Quiere seguir operando su negocio o prefiere cerrarlo ordenadamente?}
  {Es la pregunta que separa la reorganización de la liquidación, y muchas veces el
  cliente no se la ha formulado.}

  \pregunta{¿Quiénes más saben de esta situación: socios, familia, trabajadores?}
  {El procedimiento es público y se publica en el \boletin{}. El cliente debe saberlo
  antes, no después.}
\end{cuestionario}

%==============================================================================
\bloqueDocumental
%==============================================================================

Concluida la entrevista, y antes de emitir opinión alguna sobre la vía procedente, deben
reunirse los antecedentes que permiten contrastar el relato. El siguiente listado
constituye el mínimo profesional.

\begin{checklist}[Antecedentes comunes a todo procedimiento]
  \item Cédula de identidad del deudor o de su representante legal.
  \item Carpeta tributaria electrónica completa del \sii{} (para solicitar créditos, no
  la versión abreviada).
  \item Informe consolidado de deudas conforme a la normativa de la
  \cmf{}.\verificar{Denominación oficial vigente del informe y vía de obtención.}
  \item Certificado de boletín comercial y de protestos.
  \item Certificado de anotaciones vigentes del Registro Civil por cada vehículo.
  \item Certificado de dominio vigente, de hipotecas y gravámenes, y de prohibiciones,
  por cada inmueble.
  \item Cartola de la cuenta de capitalización individual y certificado de saldos de
  ahorro previsional voluntario.
  \item Cartolas bancarias de los últimos doce meses.
  \item Certificado de deuda previsional (Previred o equivalente).
\end{checklist}

\begin{checklist}[Adicionales para Empresa Deudora]
  \item Escritura de constitución, estatutos y sus modificaciones.
  \item Certificado de vigencia de la sociedad y de personería del representante.
  \item Balances y estados de resultados de los tres últimos ejercicios.
  \item Libro mayor y auxiliares de proveedores y clientes.
  \item Nómina de trabajadores con antigüedad, remuneración y estado de cotizaciones.
  \item Contratos de suministro, arrendamiento y \emph{leasing} vigentes.
  \item Detalle de existencias y activo fijo con su valorización.
\end{checklist}

\begin{foco}[Qué buscar en la carpeta tributaria]
No basta con acompañarla. Debe leerse buscando: (i)~la fecha de inicio de actividades y
la categoría, que califican al sujeto; (ii)~el nivel de ventas de los últimos ejercicios,
que define el umbral de la \Lpyme{}; (iii)~la existencia de deuda fiscal; y (iv)~las
últimas declaraciones presentadas, cuya ausencia anticipa objeciones.
\end{foco}

%==============================================================================
\bloqueFocos
%==============================================================================

\subsection{El contraste declarado / documentado}

El primer trabajo tras reunir los antecedentes consiste en construir una tabla de dos
columnas: lo que el cliente declaró y lo que los documentos acreditan. Las diferencias
se clasifican en tres categorías:

\begin{description}[leftmargin=0pt,style=nextline,font=\sffamily\bfseries]
  \item[Diferencias de memoria] Montos aproximados, fechas imprecisas. Se corrigen sin
  más.

  \item[Deudas omitidas] Obligaciones que el cliente no mencionó. Deben incorporarse y,
  si son significativas, obligan a rehacer el mapa de mayorías.

  \item[Activos o actos omitidos] Es la categoría crítica. Exige una segunda conversación
  con el cliente, documentada, antes de continuar.
\end{description}

\subsection{La decisión de vía}

Con la información verificada, la elección procedimental se reduce a un conjunto acotado
de criterios que operan en cadena. La sección \ref{sec:cap14-arbol} los sistematiza.

\subsection{El calendario real}

Un procedimiento concursal se mide en meses. El cliente que enfrenta un remate en tres
semanas necesita saber si la presentación alcanza a producir el efecto suspensivo antes
de esa fecha. Este cálculo debe hacerse en la primera semana, no cuando el remate ya
ocurrió.

%==============================================================================
\section{Árbol de decisión procedimental}
\label{sec:cap14-arbol}
%==============================================================================

\begin{flujograma}{Elección de la vía procedimental tras el diagnóstico}{cap14-arbol}
  \node[hito] (inicio) {DIAGNÓSTICO COMPLETADO\\[2pt]
    \normalfont Antecedentes verificados documentalmente};

  \coordinate (s1) at ($(inicio.south)+(0,-6mm)$);
  \node[nodo, below left=12mm and 6mm of inicio.south] (persona)
    {PERSONA DEUDORA\\[2pt]\normalfont Sin ingresos de primera categoría};
  \node[nodo, below right=12mm and 6mm of inicio.south] (empresa)
    {EMPRESA DEUDORA\\[2pt]\normalfont Contribuyente de primera categoría};

  \node[nodo, text width=44mm, below=9mm of persona] (pcrisis)
    {¿Admisible ante la \superir{}? Deudas vencidas, umbral y ausencia de ejecuciones
     notificadas};
  \node[nodo, text width=44mm, below=9mm of empresa] (ecrisis)
    {¿La empresa es viable? Contraste flujo proyectado / pasivo exigible};

  \node[favorable, text width=40mm, below left=10mm and -4mm of pcrisis.south] (reneg)
    {RENEGOCIACIÓN ADMINISTRATIVA};
  \node[adverso, text width=40mm, below right=10mm and -4mm of pcrisis.south] (liqsimp)
    {LIQUIDACIÓN SIMPLIFICADA};

  \node[favorable, text width=40mm, below left=10mm and -4mm of ecrisis.south] (reorg)
    {REORGANIZACIÓN\\[2pt]\normalfont Ordinaria o simplificada según umbral};
  \node[adverso, text width=40mm, below right=10mm and -4mm of ecrisis.south] (liq)
    {LIQUIDACIÓN\\[2pt]\normalfont Ordinaria o simplificada según umbral};

  \draw[flecha] (inicio.south) -- (s1) -| (persona.north);
  \draw[flecha] (inicio.south) -- (s1) -| (empresa.north);
  \draw[flecha] (persona) -- (pcrisis);
  \draw[flecha] (empresa) -- (ecrisis);
  \draw[flecha] (pcrisis.south) -- ++(0,-4mm) -| (reneg.north);
  \draw[flecha] (pcrisis.south) -- ++(0,-4mm) -| (liqsimp.north);
  \draw[flecha] (ecrisis.south) -- ++(0,-4mm) -| (reorg.north);
  \draw[flecha] (ecrisis.south) -- ++(0,-4mm) -| (liq.north);
\end{flujograma}

\begin{alerta}[El árbol no es automático]
Dos criterios pueden alterar el resultado: la existencia de un remate inminente, que
puede aconsejar la vía que produzca antes el efecto suspensivo; y la actitud de los
acreedores mayoritarios, que determina si una reorganización tiene mayorías alcanzables
o es un trámite destinado al rechazo.
\end{alerta}

%==============================================================================
\bloqueErrores
%==============================================================================

\errorfrecuente{Construir la nómina de pasivos con la declaración del cliente}
{Emergen acreedores no considerados; el mapa de mayorías se desarma en la junta y la
propuesta se rechaza.}
{Construirla con el informe consolidado y el boletín comercial, y usar la declaración
sólo para detectar lo que los registros no muestran (deudas informales, familiares).}

\errorfrecuente{Omitir el ahorro previsional voluntario del inventario}
{Configura la causal del \art{169 A} \Nro 1 y puede costar el \discharge{} completo.}
{Solicitar siempre el certificado de saldos de APV y cuenta 2, aunque el cliente afirme
no tenerlos.}

\errorfrecuente{No advertir sobre el CAE desde la primera reunión}
{El cliente termina el procedimiento con la deuda que más le importaba intacta, y con la
sensación de haber sido mal asesorado.}
{Advertirlo por escrito en la primera reunión y dejar constancia en el mandato.}

\errorfrecuente{Presentar sin haber calculado el calendario contra las ejecuciones en curso}
{El remate se consuma antes de que opere el efecto suspensivo.}
{Verificar el estado de cada causa ejecutiva y fijar fecha objetivo de presentación
antes de aceptar el encargo.}

\errorfrecuente{Tratar la entrevista como un trámite administrativo delegable}
{Las tres distorsiones del relato pasan sin filtro al expediente.}
{Conducirla personalmente, con el cuestionario a la vista, y documentar las respuestas
en un acta que el cliente firme.}

%==============================================================================
\section{Acta de diagnóstico}
%==============================================================================

Toda entrevista inicial debe cerrarse con un documento firmado por el cliente que
consigne lo declarado. Su función es doble: ordena el trabajo posterior y protege al
profesional frente a la alegación de que ciertos hechos fueron informados verbalmente.

\begin{escrito}[Modelo — Acta de diagnóstico concursal]
\small
\noindent\textbf{ACTA DE DIAGNÓSTICO CONCURSAL}

\medskip
\noindent En [\textsc{ciudad}], a [\textsc{fecha}], comparece [\textsc{nombre},
\textsc{rut}], en adelante «el cliente», ante [\textsc{nombre del abogado}], y declara
bajo su responsabilidad lo siguiente:

\begin{enumerate}[leftmargin=1.6em]
  \item \textbf{Calificación.} Declara [ser / no ser] contribuyente de primera categoría
  en los últimos veinticuatro meses, con ventas anuales aproximadas de
  [\textsc{monto}] UF.

  \item \textbf{Pasivo declarado.} Declara adeudar las obligaciones individualizadas en
  el Anexo A, por un total aproximado de [\textsc{monto}].

  \item \textbf{Activo declarado.} Declara ser titular de los bienes individualizados en
  el Anexo B, incluidos instrumentos de ahorro previsional voluntario y cuentas
  bancarias.

  \item \textbf{Actos del bienio anterior.} Declara [haber / no haber] celebrado, en los
  dos años anteriores, actos de disposición a título gratuito u oneroso, pagos
  anticipados u otorgamiento de garantías sobre obligaciones preexistentes. En caso
  afirmativo, se individualizan en el Anexo C.

  \item \textbf{Deudas de tratamiento especial.} Declara [tener / no tener] obligaciones
  por crédito con aval del Estado, crédito solidario universitario o pensiones de
  alimentos. \textbf{Ha sido expresamente advertido} de que la jurisprudencia sostiene
  mayoritariamente que estas obligaciones no se extinguen por efecto del procedimiento
  concursal.

  \item \textbf{Ejecuciones en curso.} Declara [haber / no haber] sido notificado de
  demandas ejecutivas, embargos o remates, individualizados en el Anexo D.

  \item \textbf{Advertencias recibidas.} El cliente declara haber sido informado de:
  (i)~el carácter público del procedimiento y su publicación en el \boletin{};
  (ii)~el efecto del \desasimiento{} sobre la administración de sus bienes;
  (iii)~las consecuencias de omitir activos o actos, incluida la eventual denegación de
  la extinción de sus saldos insolutos conforme al \art{169 A}; y (iv)~el carácter
  estimativo de los plazos informados.
\end{enumerate}

\medskip
\noindent El cliente declara que la información proporcionada es completa y veraz, y se
obliga a informar cualquier antecedente adicional que llegue a su conocimiento.

\medskip
\noindent\hfill\rule{5cm}{0.4pt}\hspace{1cm}\rule{5cm}{0.4pt}\hfill\null

\noindent\hspace{1.5cm}{\footnotesize Cliente}\hspace{5.2cm}{\footnotesize Abogado}
\end{escrito}

\begin{practica}[Valor del acta]
Además de su función probatoria, el acta cumple una función pedagógica: obliga al cliente
a leer, en frío, las advertencias que en la conversación tendió a minimizar. Es frecuente
que sea al firmarla cuando recuerda el activo o el acto que había omitido.
\end{practica}
